% In this file you should put the actual content of the blueprint.
% It will be used both by the web and the print version.
% It should *not* include the \begin{document}
%
% If you want to split the blueprint content into several files then
% the current file can be a simple sequence of \input. Otherwise It
% can start with a \section or \chapter for instance.

% === Blueprint: Spectral Characterization of Manifolds (Connes 2008) ===
% Paste this into blueprint/src/content.tex. No preamble/document env needed.

\newcommand{\ncint}{\mathchoice{\int\!\!\!\!\!\!-\!}{\int\!\!\!\!-\!}{\int\!\!-\!}{\int\!-\!}}

\chapter{Roadmap}
This blueprint follows Connes (2008) \cite{Connes2008}. First we define a spectral triple for commutative algebras as in Connes (1996) \cite{Connes1996}.


\section{Preliminaries (Chapter 2 in \cite{Connes2008})} 
This section collects all the necessary definitions, theorems and lemmas. It will be reworked when the basic structure of the proof is clear.

\begin{definition}[Involutive Algebra]\label{def:involutive-algebra}
    Let $\mathcal{A}$ be an involutive algebra over $\mathbb{C}$.
\end{definition}

\begin{definition}[Hilbert Space]\label{def:hilbert-space}
    Let $\mathcal{H}$ be a Hilbert space.
\end{definition}


\begin{definition}[Self-Adjoint Operator]\label{def:self-adjoint-operator}
    Let $D$ be a self-adjoint operator on $\mathcal{H}$.
    \uses{def:hilbert-space}
\end{definition}

\begin{definition}[Representation of Algebra]\label{def:representation-of-algebra}
    Let $\mathcal{A}$ be a represented on $\mathcal{H}$.
    \uses{def:involutive-algebra, def:hilbert-space}
\end{definition}

\begin{definition}[Spectral Dimension]\label{def:spectral-dimension}
    A self-adjoint operator $D$ has spectral dimension $p \in \mathbb{N}$ if the $n$-th characteristic value of the resolvent of $D$ is  $O(n^{-1/p})$.
    \uses{def:self-adjoint-operator}
\end{definition}

\begin{definition}[Order One]\label{def:order-one}
    A self-adjoint operator $D$ is called of order one if for all $a,b \in \mathcal{A}$, $[[D,a],b] = 0$.
    \uses{def:involutive-algebra, def:self-adjoint-operator, def:representation-of-algebra}
\end{definition}

\begin{definition}[Absolute Commutator]\label{def:absolute-commutator}
    Let $D$ be a self-adjoint operator on $\mathcal{H}$ and let $T$ be a bounded operator on $\mathcal{H}$. The absolute commutator of $D$ and $T$ is defined as $\delta(T):=[|D|,T]$.
    \uses{def:self-adjoint-operator, def:hilbert-space}
\end{definition}

\begin{lemma}[Domain of Bounded Operators]\label{lem:domain-of-bounded-operators}
    Let $D$ be a self-adjoint operator on $\mathcal{H}$ and let $T$ be a bounded operator on $\mathcal{H}$. Then $T$ maps the domain of $|D|$ into itself, i.e. $T \textrm{Dom}(|D|) \subseteq \textrm{Dom}(|D|)$.
    \uses{def:self-adjoint-operator, def:hilbert-space}
\end{lemma}

\begin{lemma}[Boundedness of Absolute Commutator]\label{lem:boundedness-of-absolute-commutator}
    Let $D$ be a self-adjoint operator on $\mathcal{H}$ and let $T$ be a bounded operator on $\mathcal{H}$. Then the absolute commutator $\delta(T) = [|D|,T]$ is a bounded operator on $\mathcal{H}$.
    \uses{def:absolute-commutator, def:self-adjoint-operator, def:hilbert-space}
\end{lemma}

\begin{definition}[Spectral Regularity] \label{def:spectral-regularity}
    Let $\mathcal{A}$ be an involutive algebra represented on a Hilbert space $\mathcal{H}$ and let $D$ be a self-adjoint operator on $\mathcal{H}$. The triple $(\mathcal{A}, \mathcal{H}, D)$ is called regular if for all $a \in \mathcal{A}$, $a$ and $[D,a]$ are in the domain $\delta^m$ for any $m \in \mathbb{N}$. 
    \uses{def:involutive-algebra, def:self-adjoint-operator, def:absolute-commutator, def:representation-of-algebra, def:hilbert-space}
\end{definition}

\begin{definition}[Differential Form Map]\label{def:differential-form-map}
    Let $\mathcal{A}$ be an involutive algebra represented on a Hilbert space $\mathcal{H}$ and let $D$ be a self-adjoint operator on $\mathcal{H}$. The differential form map
    \[
    \pi_D: \bigotimes_{j=1}^m \mathcal{A} \to \mathcal{B}(\mathcal{H})
    \] 
    is defined as
    \[
    \pi_D(a^0 \otimes a^1 \otimes \cdots \otimes a^m) =  a^0 [D,a^1] [D,a^2] \cdots [D,a^m]
    \]
    where $a^j \in \mathcal{A}$.
    \uses{def:involutive-algebra, def:self-adjoint-operator, def:representation-of-algebra, def:hilbert-space}
\end{definition}

\begin{definition}[Spectral Orientability]\label{def:spectral-orientability}
    Let $\mathcal{A}$ be an involutive algebra represented on a Hilbert space $\mathcal{H}$ and let $D$ be a self-adjoint operator on $\mathcal{H}$. Assume that $D$ has spectral dimension $p$.  

    The triple $(\mathcal{A}, \mathcal{H}, D)$ is called orientable if there exists a Hochschild cycle $c \in Z_p(\mathcal{A},\mathcal{A})$ such that $\pi_D(c)=1$ if $p$ is odd and $\pi_D(c) = \gamma$ if $p$ is even, where $\gamma$ satisfies 
    \[
    \gamma = \gamma^*, \quad \gamma^2 = 1, \quad \gamma \, D = -D \, \gamma.
    \]
    \uses{def:involutive-algebra, def:self-adjoint-operator, def:differential-form-map, def:representation-of-algebra, def:hilbert-space}  
\end{definition}

\begin{definition}[Spectral Finiteness]\label{def:spectral-finiteness}
    Let $\mathcal{A}$ be an involutive algebra represented on a Hilbert space $\mathcal{H}$ and let $D$ be a self-adjoint operator on $\mathcal{H}$. 

    The triple $(\mathcal{A}, \mathcal{H}, D)$ is called finite if the space $\mathcal{H}_\infty := \bigcap_{m \in \mathbb{N}} \textrm{Dom}(|D|^m)$ is a finitely generated projective $\mathcal{A}$-module.
    \uses{def:involutive-algebra, def:self-adjoint-operator, def:representation-of-algebra, def:hilbert-space}
\end{definition}

\begin{definition}[Spectral Absolute Continuity]\label{def:spectral-absolute-continuity}
    Let $\mathcal{A}$ be an involutive algebra represented on a Hilbert space $\mathcal{H}$ and let $D$ be a self-adjoint operator on $\mathcal{H}$. Assume that $D$ has spectral dimension $p$ and that the triple $(\mathcal{A}, \mathcal{H}, D)$ is finite.

    The triple $(\mathcal{A}, \mathcal{H}, D)$ is called absolutely continuous if the equality
    \[
    \ncint \, a(\xi\mid\eta)\,|D|^{-p} = \langle \xi, a \eta \rangle
    \]
    defines a hermitian structure $(\cdot, \cdot)$ on $\mathcal{H}_\infty$.
    \uses{def:involutive-algebra, def:self-adjoint-operator, def:representation-of-algebra, def:hilbert-space, def:spectral-finiteness}
\end{definition}

\begin{definition}[Spectral Triple]\label{def:spectral-triple}
    A spectral triple $(\mathcal{A}, \mathcal{H}, D)$ consists of an involutive algebra $\mathcal{A}$ represented on a Hilbert space $\mathcal{H}$ and a self-adjoint operator $D$ on $\mathcal{H}$ such that
    \begin{enumerate}
        \item $D$ has spectral dimension $p$,
        \item $D$ is of order one,
        \item the triple $(\mathcal{A}, \mathcal{H}, D)$ is regular,
        \item the triple $(\mathcal{A}, \mathcal{H}, D)$ is orientable,
        \item the triple $(\mathcal{A}, \mathcal{H}, D)$ is finite and
        \item the triple $(\mathcal{A}, \mathcal{H}, D)$ is absolutely continuous.
    \end{enumerate}
    \uses{def:involutive-algebra, def:self-adjoint-operator, def:representation-of-algebra, def:hilbert-space, def:spectral-dimension, def:order-one, def:spectral-regularity, def:spectral-orientability, def:spectral-finiteness, def:spectral-absolute-continuity}
\end{definition}


\begin{lemma}[Lemma 2.0.1]\label{lem:lemma-2.0.1}
    Let the triple $(\mathcal{A}, \mathcal{H}, D)$ be a finite and absolutely continuous. Then $\mathcal{H}_\infty$ is a finite projective $\mathcal{A}$-module. It can be written as $\mathcal{H}_\infty = e \mathcal{A}^n$ with $e=e^* \in M_n(\mathcal{A})$. This defines the hermitian structure so that
    \[ 
    (\xi \mid \eta) = \sum_{j} \xi_j^*  \eta_j\in \mathcal{A} \quad \forall \xi, \eta \in e \mathcal{A}^n.
    \]
    \uses{def:involutive-algebra, def:self-adjoint-operator, def:representation-of-algebra, def:hilbert-space, def:spectral-finiteness, def:spectral-absolute-continuity}
\end{lemma}


\begin{lemma}[Lemma 2.0.2]\label{lem:lemma-2.0.2}
    Let $(\mathcal{A}, \mathcal{H}, D)$ be a spectral triple. For all $a \in \mathcal{A}$ the operators $a |D|^{-p}$ are measurable and the coefficient
    \[
    \ncint \, a |D|^{-p}
    \]
    of the logarithmic divergence of their trace is unambiguously defined.
    \uses{def:spectral-triple}
\end{lemma}

\begin{proof}
    Follows from \ref{def:spectral-orientability} and theorem 8, IV.2.$\gamma$ in \cite{Connes1994} and \cite{Figueroa2000}.
    \uses{def:spectral-triple}
\end{proof}


\begin{lemma}[Lemma 2.1]\label{lem:lemma-2.1}
    The following conditions are equivalent for $T\in \mathcal{A}''$:
    \begin{enumerate}
    \item $T\in \mathcal{A}$.
    \item $[D,T]$ is bounded and, for every integer $m\ge 0$, both $T$ and $[D,T]$ belong to $\textrm{Dom}(\delta^{m})$.
    \item For every integer $m\ge 0$, $T\in \textrm{Dom}(\delta^{m})$.
    \item $T H_\infty \subset H_\infty$.
    \end{enumerate}
    \uses{def:spectral-triple}
\end{lemma}



\section{Main Theorems}

\begin{theorem}[Reconstruction Theorem 1.1]\label{thm:reconstruction-theorem-1.1}
    Let $(\mathcal{A}, \mathcal{H}, D)$ be a spectral triple such that $\mathcal{A}$ is commutative. Assume further that
    \begin{enumerate}
        \item regularity holds for all $\mathcal{A}$-endomorphisms of $\mathcal{H}_\infty$ as defined in Definition \ref{def:spectral-regularity} and
        \item the Hochschild cycle $c \in Z_p(\mathcal{A},\mathcal{A})$ in Definition \ref{def:spectral-orientability} is antisymmetric.
    \end{enumerate}
    Then there exists a compact oriented smooth manifold $M$ such that $\mathcal{A} = C^\infty(M)$.
    \uses{def:spectral-triple}
\end{theorem}
    

\begin{theorem}[Reconstruction Theorem 1.2]\label{thm:reconstruction-theorem-1.2}
    Let $(\mathcal{A}, \mathcal{H}, D)$ be a spectral triple such that $\mathcal{A}$ is commutative. Assume further that
    \begin{enumerate}
        \item the Hochschild cycle $c \in Z_p(\mathcal{A},\mathcal{A})$ in Definition \ref{def:spectral-orientability} is antisymmetric and
        \item that the multiplicity of the action of $\mathcal{A}''$ on $\mathcal{H}$ is $2^{p/2}$.
    \end{enumerate}
    Then there exists a compact oriented smooth (spin$^c$) manifold $M$ such that $\mathcal{A} = C^\infty(M)$.
    \uses{def:spectral-triple}
\end{theorem}    



\begin{corollary}[Reconstruction Theorem 1.2.a]\label{cor:reconstruction-theorem-1.2.a}
    The self-adjoint operator $D$ in Theorem \ref{thm:reconstruction-theorem-1.2} is a Dirac-type operator on $M$.
    \uses{def:spectral-triple, thm:reconstruction-theorem-1.2}
\end{corollary}


\section{Bibliography}

\begin{thebibliography}{20}
    \bibitem{Connes1994} A. Connes, \emph{Noncmmutative Geometry}, Academic Press (1994)
    \bibitem{Connes1996} A. Connes, \emph{Gravity coupled with matter and the foundation of non-commutative geometry}, Commun. Math. Phys. 182 (1996), no. 1, 155–176. \url{https://doi.org/10.1007/BF02506388}
    \bibitem{Figueroa2000} H. Figueroa, J.M. Gracia-Bond\'ia, J. Varilly, \emph{Elements of Noncommutative Geometry}, Birkh\"auser (2000) 
    \bibitem{Connes2008} A. Connes, \emph{On the spectral characterization of manifolds}, J. Noncommut. Geom. 7 (2013), no. 1, 1–82. \url{https://doi.org/10.4171/jncg/108} 
\end{thebibliography}




% --- End skeleton ---
